\documentclass[10pt,a4paper]{amsart}
\usepackage[margin=19mm]{geometry}
\usepackage{amsmath,amssymb,mathtools}
\usepackage[scr=rsfs,cal=euler]{mathalfa}
\usepackage{xfrac}
\usepackage[unicode,hidelinks,bookmarks=false]{hyperref}
\newcommand{\ie}{i.e.\ }
\newcommand{\cf}{cf.\ }
\newcommand{\resp}{respectively\ }
\newcommand{\Subsection}{\subsection}
\providecommand{\nobreakdash}{\nobreak\hbox{-}}
\setlength{\emergencystretch}{2em}
\multlinegap0pt
\usepackage{booktabs}
\usepackage{amssymb}
\newtheorem{theorem}{Theorem}
\DeclareMathOperator{\GKdim}{GKdim}
\DeclareMathOperator{\ML}{ML}
\title{Cancellation problem via \\ locally nilpotent derivations}
\author{C\'esar F. Venegas R., Helbert J. Venegas R}
\date{Source: arXiv:2512.15590v2 (CC BY 4.0)}
\begin{document}
\maketitle

\noindent\emph{Source and licence.} Cancellation problem via \\ locally nilpotent derivations. Version arXiv:2512.15590v2 (CC BY 4.0), distributed by arXiv under the Creative Commons Attribution 4.0 International licence, http://creativecommons.org/licenses/by/4.0/, as recorded in the arXiv licence metadata for this exact version and shown on the abstract page https://arxiv.org/abs/2512.15590v2. Full licence notice: \texttt{licenses/arxiv-2512.15590-CC-BY-4.0.txt}.\par\smallskip

\noindent\textit{Editorial scope.} Counted unit: the article's theorem thm:rigidity-implies-cancellation (source0.tex lines 676--678). The article's complete original proof of it is delivered with it (source0.tex lines 680--720, one \texttt{proof} environment of 41 lines). No mathematics is written, repaired or re-derived: the statement and the proof are the article's own lines, in the article's own notation, and no retained line is substituted or re-flowed.  No further source-local context is needed: every cross-reference of every retained line resolves inside the two delivered spans.  Nothing was omitted from any retained span, so the package carries no omission record.  The retained lines cite 1 works, whose entries are transcribed verbatim from the article's own bibliography source in the e-print and delivered with short editorial labels recorded in the item's reference mapping.  No retained line contains a figure environment, an image-inclusion command or drawing code, so no image is delivered and the package declares no asset dependency.  Editorial formatting only, in addition: an AMS \texttt{amsart} preamble that restates the article's own declarations and macros used by the retained lines, namely the environments \texttt{theorem} on separate counters, together with the standard packages those lines need.  The retained environments print in this document as Theorem 1 in order of appearance, whereas the article prints the same items with the numbering of its own full document.  The complete native source of the article as distributed in the arXiv e-print for this exact version is bundled unchanged as \texttt{sources/191195/source0.tex}.
\begin{theorem}(\cite[Theorem 3.3]{BellZhang2017})\label{thm:rigidity-implies-cancellation}
Let $A$ be a finitely generated domain of finite Gelfand-Kirillov (GK) dimension over a field $\Bbbk$ of characteristic zero. If $\ML(A[t]) = A$, then $A$ is cancellative.
\end{theorem}

\begin{proof}
Let $\varphi : A[t] \to B[t]$ be an isomorphism for some $\Bbbk$-algebra $B$.  
Then $\GKdim B = \GKdim A < \infty$ by \cite[Lemma 2.1(2)]{BellZhang2017}.

Let $\partial_t := \frac{\partial}{\partial t}$ be the derivation on $B[t]$; it is locally nilpotent, and the intersection of its kernels is exactly $B$.  
In particular,
\[
    \ML(B[t]) \subseteq B.
\]

On the other hand, by hypothesis,
\[
    \ML(A[t]) = A.
\]

Now let $\delta$ be a locally nilpotent derivation of $B[t]$.  
Then the derivation
\[
    \varphi^{-1} \circ \delta \circ \varphi
\]
is a locally nilpotent derivation of $A[t]$.  
Similarly, if $\delta'$ is a locally nilpotent derivation of $A[t]$, then
\[
    \varphi \circ \delta' \circ \varphi^{-1}
\]
is a locally nilpotent derivation of $B[t]$.

This shows that the locally nilpotent derivations of $A[t]$ and of $B[t]$ correspond under the isomorphism $\varphi$.  
Consequently, $\varphi$ induces an algebra isomorphism
\[
    \ML(A[t]) \;\cong\; \ML(B[t]).
\]
In particular, $\varphi$ maps $A$ onto $B$.
Consider the polynomial extension $Y = A[t]$ where $t$ has degree $1$, so that 
$Y_0 = A$ is the degree-zero homogeneous component. By 
\cite[Lemma 3.2]{BellZhang2017}, any graded isomorphism $\varphi: Y \to Y'$ 
between polynomial extensions maps degree-zero components isomorphically. 
Since $\ML(A[t]) = A$ and $\varphi$ preserves the ML-invariant (as shown above), 
we have $\ML(B[t]) = B$. Therefore $\varphi$ induces an isomorphism 
$A \cong B$, as desired.
\end{proof}

\begin{thebibliography}{99}
\bibitem[BellZh]{BellZhang2017} J.~P. Bell and J.~J. Zhang, \emph{Zariski cancellation problem for noncommutative algebras}, Selecta Math. (N.S.) 23 (2017), no.~3, 1709--1737. %
\end{thebibliography}
\end{document}
